% Algebra 2 Lesson 7-3 fast-track packet: Graphing Sine and Cosine Functions (enVision A2, Example 3 and Example 4 only)
% Ground truth: a2_7-3_TE.tex (2024 Teacher Edition transcription) and a2_T7_assessment_TE.tex. Re-cut 2026-10-31 toward the Topic 7 test:
% topic items 17 (Example 3), 11 and 18 (Example 4). Quiz items taught: 1 (Example 3), 2 (Example 4), 4 (Example 3B, Try It 4b, extension Part E).
% Shared preamble for the "fast" (hacked-down) Algebra 2 lesson packets, Oct 2026.
% Each packet: % Shared preamble for the "fast" (hacked-down) Algebra 2 lesson packets, Oct 2026.
% Each packet: % Shared preamble for the "fast" (hacked-down) Algebra 2 lesson packets, Oct 2026.
% Each packet: \input{fast-packet-template} then \begin{document} ... \end{document}.
\documentclass[12pt]{article}
\usepackage[margin=0.7in]{geometry}
\usepackage{amsmath,amssymb}
\usepackage{booktabs,tabularx,array}
\usepackage{enumitem}
\usepackage[T1]{fontenc}
\usepackage[utf8]{inputenc}
\usepackage{lmodern}
\usepackage{tikz}
\usepackage{pgfplots}
\pgfplotsset{compat=1.18}
\usepackage{xcolor}
\setlist{nosep,leftmargin=1.4em}
\setlength{\parindent}{0pt}
\setlength{\parskip}{6pt}
\renewcommand{\arraystretch}{1.35}
\newcolumntype{Y}{>{\raggedright\arraybackslash}X}
\pagestyle{empty}

% blanks and workspace
\newcommand{\blank}[1][2.6cm]{\rule[-2pt]{#1}{0.5pt}}
\newcommand{\lines}[1][1]{\par\vspace{2pt}\foreach \i in {1,...,#1}{\rule{0pt}{22pt}\rule{\linewidth}{0.4pt}\par}}
\newcommand{\wspart}[2]{\vspace{4pt}{\large\bfseries #1: #2}\par\vspace{-6pt}\rule{\textwidth}{0.4pt}\par}

% header: \packethead{Lesson 1-3}{Piecewise-Defined Functions}
\newcommand{\packethead}[2]{%
{\LARGE\bfseries #2}\hfill Name \blank[5cm]\\[2pt]
Algebra 2 \quad #1 (fast track) \hfill Date \blank[2.5cm] \quad Period \blank[1.2cm]\par\vspace{2pt}}

% boxed goals: \goals{learning goal}{language goal}
\newcommand{\goals}[2]{%
\fbox{\parbox{\dimexpr\textwidth-2\fboxsep-2\fboxrule}{%
\textbf{Learning goal:} #1\\[2pt]
\textbf{Language goal:} #2}}\par}

% quiz target box: \quiztarget{text}
\newcommand{\quiztarget}[1]{%
\vspace{4pt}\colorbox{gray!12}{\parbox{\dimexpr\textwidth-2\fboxsep}{\textbf{On the quiz:} #1}}\par}

% IXL table row helper: code & skill & which quiz question it feeds
\newcommand{\ixlrow}[3]{\texttt{#1} & #2 & #3 \\}

% one grid style for every graph; override xmin/xmax/... per plot
\pgfplotsset{sheetgrid/.style={
  axis lines=middle, grid=both, grid style={gray!45},
  xmin=-10.5, xmax=10.5, ymin=-10.5, ymax=10.5,
  xtick={-10,-8,...,10}, ytick={-10,-8,...,10},
  x=0.28cm, y=0.28cm,
  ticklabel style={font=\scriptsize, fill=white, inner sep=1pt},
  xlabel={$x$}, ylabel={$y$},
  axis line style={thick}, clip=true}}
% unit grid for small windows
\pgfplotsset{unitgrid/.style={sheetgrid,
  xmin=-5.5, xmax=5.5, ymin=-5.5, ymax=5.5,
  xtick={-5,...,5}, ytick={-5,...,5}, x=0.5cm, y=0.5cm}}
% open and closed endpoints
\tikzset{closed/.style={circle, fill, inner sep=1.6pt},
         open/.style={circle, draw, fill=white, inner sep=1.4pt}}
 then \begin{document} ... \end{document}.
\documentclass[12pt]{article}
\usepackage[margin=0.7in]{geometry}
\usepackage{amsmath,amssymb}
\usepackage{booktabs,tabularx,array}
\usepackage{enumitem}
\usepackage[T1]{fontenc}
\usepackage[utf8]{inputenc}
\usepackage{lmodern}
\usepackage{tikz}
\usepackage{pgfplots}
\pgfplotsset{compat=1.18}
\usepackage{xcolor}
\setlist{nosep,leftmargin=1.4em}
\setlength{\parindent}{0pt}
\setlength{\parskip}{6pt}
\renewcommand{\arraystretch}{1.35}
\newcolumntype{Y}{>{\raggedright\arraybackslash}X}
\pagestyle{empty}

% blanks and workspace
\newcommand{\blank}[1][2.6cm]{\rule[-2pt]{#1}{0.5pt}}
\newcommand{\lines}[1][1]{\par\vspace{2pt}\foreach \i in {1,...,#1}{\rule{0pt}{22pt}\rule{\linewidth}{0.4pt}\par}}
\newcommand{\wspart}[2]{\vspace{4pt}{\large\bfseries #1: #2}\par\vspace{-6pt}\rule{\textwidth}{0.4pt}\par}

% header: \packethead{Lesson 1-3}{Piecewise-Defined Functions}
\newcommand{\packethead}[2]{%
{\LARGE\bfseries #2}\hfill Name \blank[5cm]\\[2pt]
Algebra 2 \quad #1 (fast track) \hfill Date \blank[2.5cm] \quad Period \blank[1.2cm]\par\vspace{2pt}}

% boxed goals: \goals{learning goal}{language goal}
\newcommand{\goals}[2]{%
\fbox{\parbox{\dimexpr\textwidth-2\fboxsep-2\fboxrule}{%
\textbf{Learning goal:} #1\\[2pt]
\textbf{Language goal:} #2}}\par}

% quiz target box: \quiztarget{text}
\newcommand{\quiztarget}[1]{%
\vspace{4pt}\colorbox{gray!12}{\parbox{\dimexpr\textwidth-2\fboxsep}{\textbf{On the quiz:} #1}}\par}

% IXL table row helper: code & skill & which quiz question it feeds
\newcommand{\ixlrow}[3]{\texttt{#1} & #2 & #3 \\}

% one grid style for every graph; override xmin/xmax/... per plot
\pgfplotsset{sheetgrid/.style={
  axis lines=middle, grid=both, grid style={gray!45},
  xmin=-10.5, xmax=10.5, ymin=-10.5, ymax=10.5,
  xtick={-10,-8,...,10}, ytick={-10,-8,...,10},
  x=0.28cm, y=0.28cm,
  ticklabel style={font=\scriptsize, fill=white, inner sep=1pt},
  xlabel={$x$}, ylabel={$y$},
  axis line style={thick}, clip=true}}
% unit grid for small windows
\pgfplotsset{unitgrid/.style={sheetgrid,
  xmin=-5.5, xmax=5.5, ymin=-5.5, ymax=5.5,
  xtick={-5,...,5}, ytick={-5,...,5}, x=0.5cm, y=0.5cm}}
% open and closed endpoints
\tikzset{closed/.style={circle, fill, inner sep=1.6pt},
         open/.style={circle, draw, fill=white, inner sep=1.4pt}}
 then \begin{document} ... \end{document}.
\documentclass[12pt]{article}
\usepackage[margin=0.7in]{geometry}
\usepackage{amsmath,amssymb}
\usepackage{booktabs,tabularx,array}
\usepackage{enumitem}
\usepackage[T1]{fontenc}
\usepackage[utf8]{inputenc}
\usepackage{lmodern}
\usepackage{tikz}
\usepackage{pgfplots}
\pgfplotsset{compat=1.18}
\usepackage{xcolor}
\setlist{nosep,leftmargin=1.4em}
\setlength{\parindent}{0pt}
\setlength{\parskip}{6pt}
\renewcommand{\arraystretch}{1.35}
\newcolumntype{Y}{>{\raggedright\arraybackslash}X}
\pagestyle{empty}

% blanks and workspace
\newcommand{\blank}[1][2.6cm]{\rule[-2pt]{#1}{0.5pt}}
\newcommand{\lines}[1][1]{\par\vspace{2pt}\foreach \i in {1,...,#1}{\rule{0pt}{22pt}\rule{\linewidth}{0.4pt}\par}}
\newcommand{\wspart}[2]{\vspace{4pt}{\large\bfseries #1: #2}\par\vspace{-6pt}\rule{\textwidth}{0.4pt}\par}

% header: \packethead{Lesson 1-3}{Piecewise-Defined Functions}
\newcommand{\packethead}[2]{%
{\LARGE\bfseries #2}\hfill Name \blank[5cm]\\[2pt]
Algebra 2 \quad #1 (fast track) \hfill Date \blank[2.5cm] \quad Period \blank[1.2cm]\par\vspace{2pt}}

% boxed goals: \goals{learning goal}{language goal}
\newcommand{\goals}[2]{%
\fbox{\parbox{\dimexpr\textwidth-2\fboxsep-2\fboxrule}{%
\textbf{Learning goal:} #1\\[2pt]
\textbf{Language goal:} #2}}\par}

% quiz target box: \quiztarget{text}
\newcommand{\quiztarget}[1]{%
\vspace{4pt}\colorbox{gray!12}{\parbox{\dimexpr\textwidth-2\fboxsep}{\textbf{On the quiz:} #1}}\par}

% IXL table row helper: code & skill & which quiz question it feeds
\newcommand{\ixlrow}[3]{\texttt{#1} & #2 & #3 \\}

% one grid style for every graph; override xmin/xmax/... per plot
\pgfplotsset{sheetgrid/.style={
  axis lines=middle, grid=both, grid style={gray!45},
  xmin=-10.5, xmax=10.5, ymin=-10.5, ymax=10.5,
  xtick={-10,-8,...,10}, ytick={-10,-8,...,10},
  x=0.28cm, y=0.28cm,
  ticklabel style={font=\scriptsize, fill=white, inner sep=1pt},
  xlabel={$x$}, ylabel={$y$},
  axis line style={thick}, clip=true}}
% unit grid for small windows
\pgfplotsset{unitgrid/.style={sheetgrid,
  xmin=-5.5, xmax=5.5, ymin=-5.5, ymax=5.5,
  xtick={-5,...,5}, ytick={-5,...,5}, x=0.5cm, y=0.5cm}}
% open and closed endpoints
\tikzset{closed/.style={circle, fill, inner sep=1.6pt},
         open/.style={circle, draw, fill=white, inner sep=1.4pt}}

% radian ticks: multiples of pi with labels, pi/4 or pi/2 minor grid
\pgfplotsset{trigx/.style={
  xtick={-12.566,-9.4248,-6.2832,-3.1416,3.1416,6.2832,9.4248,12.566},
  xticklabels={$-4\pi$,$-3\pi$,$-2\pi$,$-\pi$,$\pi$,$2\pi$,$3\pi$,$4\pi$},
  minor grid style={gray!30}}}
\begin{document}
\relpenalty=10000 \binoppenalty=10000 % keep each inline formula on one line

\packethead{Lesson 7-3}{Graphing Sine and Cosine Functions}
\goals{I can find the amplitude, period and frequency of a sine or cosine function from its equation and graph it from five key points.}{I can explain how $a$ and $b$ change the graph of $y=a\sin bx$ using the words \emph{amplitude}, \emph{period}, \emph{midline}, and \emph{frequency}.}

\quiztarget{(1) Amplitude and period: Example 3. (2) Rate of change from key points: Example 4, Practice. (3) Equation against a graph (full Savvas quiz only). (4) Reflect $5\cos(2x)$: Example 3B, Try It 4b. (5) Tide (full Savvas quiz only).}
\topictarget{7}{Question 17 asks for a frequency (Example 3); 11 and 18 ask for an average rate of change (Example 4). Questions 2 (equation against a graph), 6 (tide) and 12 (repeating sine values) come from Examples 6, 5 and 1, which this packet skips.}

\wspart{Example 3}{Identify Amplitude and Period}
% Rule from the TE Reteach box (a != 0, b > 0: period 2pi/b); written with |b| so it stays true for sin(-2x) in Practice 23.
For $y=a\sin bx$ or $y=a\cos bx$ ($a\neq0$, $b\neq0$): the \textbf{amplitude} is $|a|$ (the distance from the \textbf{midline} to a maximum or minimum) and the \textbf{period} is $\frac{2\pi}{|b|}$ (the length of one cycle).

\begin{minipage}[t]{0.47\textwidth}\raggedright
\textbf{A.} What are the amplitude and period of $f(x)=3\cos x$?
\begin{enumerate}[label=\arabic*.]
\item Graph $f(x)=3\cos x$ and the parent $y=\cos x$.
\item Midline $y=0$. Maximum values 3, minimum values $-3$: 3 times those of $\cos x$.
\item The amplitude of $f(x)=3\cos x$ is 3, and the period is $2\pi$.
\end{enumerate}
\begin{center}
\begin{tikzpicture}
\begin{axis}[sheetgrid, trigx, xmin=-11.2, xmax=11.2, ymin=-5.5, ymax=5.5,
  xtick={-9.4248,-6.2832,-3.1416,3.1416,6.2832,9.4248},
  xticklabels={{},$-2\pi$,{},{},$2\pi$,{}},
  minor xtick={-10.9956,-9.4248,...,11}, ytick={-4,-2,2,4}, minor ytick={-5,...,5},
  x=0.33cm, y=0.2cm, ticklabel style={font=\tiny, fill=white, inner sep=0.5pt}, set layers, axis on top, grid style={gray!45, /pgfplots/on layer=axis background}, minor grid style={gray!30, /pgfplots/on layer=axis background},
  title={\scriptsize\textcolor{red}{$f(x)=3\cos x$} \quad \textcolor{blue}{$y=\cos x$}}, title style={yshift=-4pt}]
\addplot[very thick, blue, domain=-10.9:10.9, samples=200, <->] {cos(deg(x))};
\addplot[very thick, red, domain=-10.9:10.9, samples=200, <->] {3*cos(deg(x))};
\addplot[thick, dashed, red] coordinates {(6.2832,0) (6.2832,3)};
\addplot[thick, dashed, red] coordinates {(-6.2832,3) (0,3)};
\addplot[thick, dashed, blue] coordinates {(-6.2832,1) (0,1)};
\node[font=\tiny, red, fill=white, inner sep=0.5pt, anchor=west] at (axis cs:6.6,1.5) {3};
\node[font=\tiny, inner sep=1pt, anchor=north east] at (axis cs:0,0) {$O$};
% odd multiples of pi: labels above the axis so they do not sit on the minima of either curve
\node[font=\tiny, fill=white, inner sep=0.5pt, anchor=south, yshift=1pt] at (axis cs:-9.4248,0) {$-3\pi$};
\node[font=\tiny, fill=white, inner sep=0.5pt, anchor=south, yshift=1pt] at (axis cs:-3.1416,0) {$-\pi$};
\node[font=\tiny, fill=white, inner sep=0.5pt, anchor=south, yshift=1pt] at (axis cs:3.1416,0) {$\pi$};
\node[font=\tiny, fill=white, inner sep=0.5pt, anchor=south, yshift=1pt] at (axis cs:9.4248,0) {$3\pi$};
\end{axis}
\end{tikzpicture}
\end{center}
\vspace{-6pt}\textbf{Amplitude 3, period $2\pi$.}
\end{minipage}\hfill
\begin{minipage}[t]{0.50\textwidth}\raggedright
\textbf{B.} What are the amplitude and period of $y=-\sin 2x$?
\begin{enumerate}[label=\arabic*.]
\item $y=-f(x)$ is a \textbf{reflection} of the parent $y=\sin x$ over the $x$-axis.
\item $y=f(2x)$ is a horizontal compression by a factor of $\frac12$.
\item Neither one changes the amplitude: $|-1|=1$.
\item Period compressed by $\frac12$: $\frac12(2\pi)=\pi$.
\item Check with a graph:
\end{enumerate}
\begin{center}
\begin{tikzpicture}
\begin{axis}[sheetgrid, trigx, xmin=-8.2, xmax=8.2, ymin=-1.35, ymax=1.35,
  xtick={-6.2832,-3.1416,3.1416,6.2832}, xticklabels={$-2\pi$,$-\pi$,$\pi$,$2\pi$},
  minor xtick={-7.854,-6.2832,...,8}, ytick={-1,1}, yticklabels={$-1$,{}}, minor ytick={-1.25,-1,...,1.3},
  x=0.45cm, y=0.55cm, ticklabel style={font=\tiny, fill=white, inner sep=0.5pt}, xticklabel style={anchor=north east}, set layers, axis on top, grid style={gray!45, /pgfplots/on layer=axis background}, minor grid style={gray!30, /pgfplots/on layer=axis background},
  title={\scriptsize\textcolor{red}{$y=-\sin 2x$}}, title style={yshift=-4pt}]
\addplot[very thick, red, domain=-7.7:7.7, samples=240, <->] {-sin(deg(2*x))};
\addplot[very thick, green!60!black] coordinates {(-3.927,0) (-3.927,1)};
\node[font=\tiny, inner sep=1pt, anchor=north east] at (axis cs:0,0) {$O$};
\end{axis}
\end{tikzpicture}
\end{center}
\vspace{-6pt}Repeats every $\pi$. \textbf{Amplitude 1, period $\pi$.}
\end{minipage}

\textbf{Say it:} The amplitude of $y=-\sin 2x$ is \blank[1.2cm] because the midline is \blank[1.5cm] and a maximum is \blank[1.2cm] above it. The period is \blank[1.2cm] because \lines[1]

\wspart{Try It 3}{What are the amplitude and period of each function?}
\begin{minipage}[t]{0.48\textwidth}
\textbf{a.} $y=\frac13\cos\frac12 x$\\[2pt]
amplitude: \blank[1.8cm] \quad period: \blank[1.8cm]
\end{minipage}\hfill
\begin{minipage}[t]{0.48\textwidth}
\textbf{b.} $y=2\sin\pi x$\\[2pt]
amplitude: \blank[1.8cm] \quad period: \blank[1.8cm]
\end{minipage}
\newpage

% TE Frequency concept box (between Try It 3 and Example 4). Topic 7 test item 17 asks for a frequency.
\textbf{Frequency} is the reciprocal of the period: the number of complete cycles in one unit of $x$. $y=\sin x$ and $y=\cos x$ have period $2\pi$ and frequency $\frac{1}{2\pi}$; $y=-\sin 2x$ has period $\pi$ and frequency $\frac{1}{\pi}$.

\wspart{Example 4}{Graph $y=a\sin bx$ and $y=a\cos bx$}
\begin{minipage}[t]{0.49\textwidth}\raggedright
\textbf{A.} Graph $f(x)=5\sin\frac12 x$.
\begin{enumerate}[label=\arabic*.]
\item $a=5$: vertical stretch, so the amplitude is 5.
\item $b=\frac12$: horizontal stretch by a factor of 2, so the period is $\frac{2\pi}{1/2}=4\pi$.
\item Divide one period into fourths: $\frac{4\pi}{4}=\pi$. Plot, in order, a zero, a maximum, a zero, a minimum and a zero.
\end{enumerate}
\begin{center}
\begin{tikzpicture}
\begin{axis}[sheetgrid, xmin=-1.9, xmax=20.9, ymin=-6.5, ymax=6.5,
  xtick={6.2832,12.566,18.850}, xticklabels={{},{},{}}, minor xtick={-1.5708,0,...,20.5},
  ytick={-4,-2,2,4}, minor ytick={-6,...,6}, x=0.31cm, y=0.27cm,
  ticklabel style={font=\tiny, fill=white, inner sep=0.5pt}, set layers, axis on top, grid style={gray!45, /pgfplots/on layer=axis background}, minor grid style={gray!30, /pgfplots/on layer=axis background},
  title={\scriptsize\textcolor{green!50!black}{$f(x)=5\sin\frac12x$}}, title style={yshift=-4pt}]
\addplot[very thick, green!50!black, domain=-1.5:20.3, samples=200, <->] {5*sin(deg(x/2))};
\node[closed, red] at (axis cs:0,0) {};
\node[closed, red] at (axis cs:3.1416,5) {};
\node[closed, red] at (axis cs:6.2832,0) {};
\node[closed, red] at (axis cs:9.4248,-5) {};
\node[closed, red] at (axis cs:12.566,0) {};
\node[font=\tiny, fill=white, inner sep=0.5pt, anchor=north west] at (axis cs:0.2,-0.2) {$(0,0)$};
\node[font=\tiny, fill=white, inner sep=0.5pt, anchor=south] at (axis cs:3.1416,5.3) {$(\pi,5)$};
\node[font=\tiny, fill=white, inner sep=0.5pt, anchor=south west] at (axis cs:6.4,0.2) {$(2\pi,0)$};
\node[font=\tiny, fill=white, inner sep=0.5pt, anchor=north] at (axis cs:9.4248,-5.3) {$(3\pi,-5)$};
\node[font=\tiny, fill=white, inner sep=0.5pt, anchor=south east] at (axis cs:12.4,0.2) {$(4\pi,0)$};
\node[font=\tiny, fill=white, inner sep=0.5pt, anchor=north east] at (axis cs:6.2,-0.1) {$2\pi$};
\node[font=\tiny, fill=white, inner sep=0.5pt, anchor=north west] at (axis cs:12.7,-0.1) {$4\pi$};
\node[font=\tiny, fill=white, inner sep=0.5pt, anchor=north east] at (axis cs:18.7,-0.1) {$6\pi$};
\end{axis}
\end{tikzpicture}
\end{center}
\vspace{-6pt}Domain $(-\infty,\infty)$, range $[-5,5]$. \textbf{Amplitude 5, period $4\pi$, frequency $\frac{1}{4\pi}$.}
\end{minipage}\hfill
\begin{minipage}[t]{0.49\textwidth}\raggedright
\textbf{B.} Graph $f(x)=3\cos\pi x$.
\begin{enumerate}[label=\arabic*.]
\item $a=3$: vertical stretch by a factor of 3, so the amplitude is 3.
\item $b=\pi$: horizontal compression by a factor of $\frac1\pi$, so the period is $\frac{2\pi}{\pi}=2$.
\item Fourths: $\frac{2}{4}=\frac12$. Plot, in order, a maximum, a zero, a minimum, a zero and a maximum.
\end{enumerate}
\begin{center}
\begin{tikzpicture}
\begin{axis}[sheetgrid, xmin=-0.75, xmax=5.25, ymin=-4.5, ymax=4.5,
  xtick={1,2,3,4}, minor xtick={-0.5,0,...,5}, ytick={-2,2}, minor ytick={-4,...,4},
  x=1.2cm, y=0.36cm, ticklabel style={font=\tiny, fill=white, inner sep=0.5pt}, xticklabel style={anchor=north west}, set layers, axis on top, grid style={gray!45, /pgfplots/on layer=axis background}, minor grid style={gray!30, /pgfplots/on layer=axis background},
  title={\scriptsize\textcolor{red}{$f(x)=3\cos\pi x$}}, title style={yshift=-4pt}]
\addplot[very thick, red, domain=-0.6:5.1, samples=200, <->] {3*cos(deg(pi*x))};
\node[closed, red] at (axis cs:0,3) {};
\node[closed, red] at (axis cs:0.5,0) {};
\node[closed, red] at (axis cs:1,-3) {};
\node[closed, red] at (axis cs:1.5,0) {};
\node[closed, red] at (axis cs:2,3) {};
\node[font=\tiny, fill=white, inner sep=0.5pt, anchor=south east] at (axis cs:-0.04,3.15) {$(0,3)$};
\node[font=\tiny, fill=white, inner sep=0.5pt, anchor=south west] at (axis cs:0.55,0.15) {$(\frac12,0)$};
\node[font=\tiny, fill=white, inner sep=0.5pt, anchor=north] at (axis cs:1,-3.3) {$(1,-3)$};
\node[font=\tiny, fill=white, inner sep=0.5pt, anchor=south west] at (axis cs:1.6,0.1) {$(\frac32,0)$};
\node[font=\tiny, fill=white, inner sep=0.5pt, anchor=south] at (axis cs:2,3.3) {$(2,3)$};
\end{axis}
\end{tikzpicture}
\end{center}
\vspace{-6pt}Domain $(-\infty,\infty)$, range $[-3,3]$. \textbf{Amplitude 3, period 2, frequency $\frac12$.}
\end{minipage}

\vspace{2pt}
{\small When $|b|<1$ the graph is stretched horizontally and the period is longer than $2\pi$; when $|b|>1$ it is compressed and the period is shorter.}

\textbf{Say it:} In $f(x)=5\sin\frac12 x$, the amplitude is \blank[1.2cm] because $a=$ \blank[1cm]. The period is \blank[1.2cm] because $2\pi\div$ \blank[1cm] $=$ \blank[1cm], so the frequency is \blank[1.2cm] and the key points are \blank[1.2cm] apart, starting on the \blank[1.6cm] $y=0$.

\wspart{Try It 4}{Graph the functions. What are the key features?}
\begin{minipage}[t]{0.49\textwidth}\raggedright
\textbf{a.} $g(x)=\frac32\cos\frac34 x$
\begin{center}
\begin{tikzpicture}
\begin{axis}[sheetgrid, xmin=-4.5, xmax=14.5, ymin=-3.4, ymax=3.4,
  xtick={4,8,12}, minor xtick={-4,...,14}, ytick={-2,2}, minor ytick={-3,...,3},
  x=0.4cm, y=0.5cm, ticklabel style={font=\tiny, fill=white, inner sep=0.5pt}]
\node[font=\tiny, inner sep=1pt, anchor=north east] at (axis cs:0,0) {$O$};
\end{axis}
\end{tikzpicture}
\end{center}
amplitude \blank[1.4cm] \; period \blank[1.4cm]\\[4pt]
frequency \blank[1.4cm] \; range \blank[1.6cm]
\end{minipage}\hfill
\begin{minipage}[t]{0.49\textwidth}\raggedright
\textbf{b.} $h(x)=-\sin 2\pi x$
\begin{center}
\begin{tikzpicture}
\begin{axis}[sheetgrid, xmin=-1.25, xmax=3.75, ymin=-1.75, ymax=1.75,
  xtick={1,2,3}, minor xtick={-1,-0.5,...,3.6}, ytick={-1,1}, minor ytick={-1.5,-1,...,1.6},
  x=1.5cm, y=1cm, ticklabel style={font=\tiny, fill=white, inner sep=0.5pt}]
\node[font=\tiny, inner sep=1pt, anchor=north east] at (axis cs:0,0) {$O$};
\end{axis}
\end{tikzpicture}
\end{center}
amplitude \blank[1.4cm] \; period \blank[1.4cm]\\[4pt]
frequency \blank[1.4cm] \; range \blank[1.6cm]
\end{minipage}

\newpage

% R5: structural parallel of quiz item 2 (average rate of change of y = a sin bx from 0 to the
% quarter-period key point); no Try It rehearses it. Topic 7 test items 11 and 18 use the same skill.
\wspart{Practice}{Average rate of change from the key points}
Use the key points you plotted in Example 4.

{\small Hint: the average rate of change of $f$ on the interval $[p,q]$ is $\dfrac{f(q)-f(p)}{q-p}$.}

\begin{minipage}[t]{0.48\textwidth}
\textbf{a.} What is the average rate of change of $f(x)=5\sin\frac12 x$ on the interval $[0,\pi]$?\lines[2]
\end{minipage}\hfill
\begin{minipage}[t]{0.48\textwidth}
\textbf{b.} What is the average rate of change of $f(x)=3\cos\pi x$ on the interval $[0,1]$?\lines[2]
\end{minipage}

% R3 DOK 3: Examples 3 and 4 carry no DOK 3 item (TE item analysis), so the lesson's DOK 3
% Practice & Problem Solving 23 (filed under the cut Example 6, Performance Task) is the named
% self-paced extension, verbatim. Part E rehearses bonus quiz item 4 (reflection over the x-axis).
\wspart{Extension}{DOK 3, Performance Task (Practice 23)}
{\small Self-paced; from Example 6, which this packet skips. Parts A and C use Example 4's five key points.}

Dominga is investigating how the signs of the parameters $a$ and $b$ create transformations of the sine function.

\begin{minipage}[t]{0.49\textwidth}\raggedright
\textbf{Part A} Graph $f(x)=\sin 2x$ and $g(x)=-\sin(2x)$ on the same coordinate plane.
\begin{center}
\begin{tikzpicture}
\begin{axis}[sheetgrid, trigx, xmin=-6.6, xmax=6.6, ymin=-1.35, ymax=1.35,
  xtick={-6.2832,-3.1416,3.1416,6.2832}, xticklabels={$-2\pi$,$-\pi$,$\pi$,$2\pi$},
  minor xtick={-6.2832,-5.4978,...,6.3}, ytick={-1,-0.5,0.5,1}, minor ytick={-1.25,-1,...,1.3},
  x=0.6cm, y=1.1cm, ticklabel style={font=\tiny, fill=white, inner sep=0.5pt}]
\node[font=\tiny, inner sep=1pt, anchor=north east] at (axis cs:0,0) {$O$};
\end{axis}
\end{tikzpicture}
\end{center}
\textbf{Part B} How are the graphs of $f(x)=\sin(2x)$ and $g(x)=-\sin(2x)$ related?\lines[2]
\end{minipage}\hfill
\begin{minipage}[t]{0.49\textwidth}\raggedright
\textbf{Part C} Graph $f(x)=\sin(2x)$ and $h(x)=\sin(-2x)$ on the same coordinate plane.
\begin{center}
\begin{tikzpicture}
\begin{axis}[sheetgrid, trigx, xmin=-6.6, xmax=6.6, ymin=-1.35, ymax=1.35,
  xtick={-6.2832,-3.1416,3.1416,6.2832}, xticklabels={$-2\pi$,$-\pi$,$\pi$,$2\pi$},
  minor xtick={-6.2832,-5.4978,...,6.3}, ytick={-1,-0.5,0.5,1}, minor ytick={-1.25,-1,...,1.3},
  x=0.6cm, y=1.1cm, ticklabel style={font=\tiny, fill=white, inner sep=0.5pt}]
\node[font=\tiny, inner sep=1pt, anchor=north east] at (axis cs:0,0) {$O$};
\end{axis}
\end{tikzpicture}
\end{center}
\textbf{Part D} How are the graphs of $f(x)=\sin 2x$ and $h(x)=\sin(-2x)$ related?\lines[2]
\end{minipage}

\vspace{2pt}
{\small Hint: for Parts A and C, one period is $\frac{2\pi}{2}=\pi$. Divide it into fourths and plot a zero, a maximum or minimum, a zero, the other one, and a zero.}

\textbf{Part E} How is the graph of $y=a\sin(bx)$ affected when $a$ or $b$ is replaced with its opposite? Explain.\lines[5]

\newpage
\wspart{IXL}{Practice for the quiz}
Do the starred skills first. \textbf{SmartScore 70 or above on each starred ($\bigstar$) skill replaces your quiz score.}

\vspace{2pt}
\begin{tabularx}{\textwidth}{l Y Y}
\toprule
\textbf{Code} & \textbf{Skill} & \textbf{Feeds} \\
\midrule
\ixlrow{6H8}{Evaluate trigonometric ratios using reference angles: sin, cos, tan}{Warm-up: the sine and cosine values at the key points (Lesson 7-2)}
\ixlrow{9S8 $\bigstar$}{Find properties of sine and cosine functions}{Quiz 1: amplitude $|a|$ and period $2\pi/|b|$ from the equation (Example 3); Quiz 4 bonus: a negative $a$ keeps the amplitude}
\ixlrow{A7V $\bigstar$}{Graph sine and cosine functions}{Quiz 2: the five key points of $y=a\sin bx$ (Example 4); Quiz 4 bonus: graph a reflection over the $x$-axis}
\ixlrow{5XN}{Write equations of sine and cosine functions from graphs}{Quiz 3 (full Savvas quiz only): amplitude and period read off a graph}
\bottomrule
\end{tabularx}

\vspace{8pt}
\textbf{Before the quiz, finish these sentences.} {\small (Spanish: amplitude = \emph{amplitud}, period = \emph{período}, midline = \emph{línea media}, frequency = \emph{frecuencia}, periodic function = \emph{función periódica}.)}

The \blank[2.5cm] is the distance from the midline to a maximum or a minimum, so the amplitude of $y=a\cos bx$ is \blank[2cm].

The \blank[2.5cm] is the length of one cycle, so the period of $y=a\sin bx$ is \blank[2cm].

The \blank[2.2cm] is the reciprocal of the period, so a function with period $4\pi$ has frequency \blank[1.8cm].

The five key points of $y=a\sin bx$ are \blank[2.2cm] of a period apart, and the first one is on the \blank[2.2cm].

Changing $a$ to $-a$ does not change the amplitude; it \blank[6cm].

\vspace{6pt}
\textbf{Exit:} Explain how $a$ and $b$ change the graph of $y=4\cos 2x$. The \blank[2.2cm] is \blank[1.4cm] because $a=4$, and the \blank[2.2cm] is \blank[1.4cm] because \lines[2]
Its endpoint values on $\left[0,\frac{\pi}{2}\right]$ are $(0,\ \blank[1cm])$ and $\left(\frac{\pi}{2},\ \blank[1cm]\right)$, so its average rate of change on $\left[0,\frac{\pi}{2}\right]$ is \blank[2cm].

\end{document}
